\section{Sequentielle quadratische Programmierung (SQP)}
\subsection{Problemstruktur}
Wir betrachten
\[
\min_{\mathbf x}J(\mathbf x)
\quad\text{unter}\quad g(\mathbf x)=0,\;h(\mathbf x)\leq0.
\]
\begin{definition}[Parametervektor]
\[
\mathbf x=(\kappa_0,\dots,\kappa_N,\phi_0,\dots,\phi_N).
\]
\end{definition}

\subsection{Residuenformulierung}
\begin{definition}[Residuenvektor]
\[
J(\mathbf x)=\frac12\|\mathbf r(\mathbf x)\|^2.
\]
\end{definition}
Typische Residuenkomponenten sind
\[
r_i^{(\kappa)}=\sqrt\gamma\frac{\kappa_{i+1}-\kappa_i}{\Delta s},
\quad
r_i^{(\phi)}=\sqrt\delta\frac{\phi_{i+1}-\phi_i}{\Delta s},
\]
\[
r_i^{(\mathrm{dyn})}=\sqrt\beta(v^2\kappa_i-g\sin\phi_i).
\]

\subsection{Gauss--Newton-Schritt}
\begin{definition}[SQP-Schritt]
In Iteration \(k\) wird \(\mathbf p_k\) berechnet aus
\[
(J_r^T J_r)\mathbf p_k=-J_r^T\mathbf r.
\]
\end{definition}
Die Hesse-Matrix wird approximiert durch \(H\approx J_r^T J_r\).

\subsection{Dünn besetzte Struktur}
Die Systemmatrix besitzt Bandstruktur: tridiagonale Kopplung für
Glattheitsterme und diagonale Beiträge dynamischer Terme.

\subsection{Behandlung der Bedingungen}
Bedingungen werden durch Projektion, Strafglieder oder Lagrange-Multiplikatoren
behandelt.

\subsection{Iterationsschema}
\begin{algorithm}
\caption{SQP-Iteration}
\begin{algorithmic}
\State initialisiere $\mathbf{x}_0$
\For{$k=0,1,\dots$}
  \State berechne $\mathbf r(\mathbf x_k)$
  \State berechne $J_r(\mathbf x_k)$
  \State löse $(J_r^T J_r)\mathbf p_k=-J_r^T\mathbf r$
  \State wähle Schrittweite $\alpha_k$
  \State aktualisiere $\mathbf x_{k+1}=\mathbf x_k+\alpha_k\mathbf p_k$
\EndFor
\end{algorithmic}
\end{algorithm}

\subsection{Interpretation}
\begin{summary}
SQP bietet ein praktisches und effizientes Verfahren zur Lösung des diskreten
Alignment-Optimierungsproblems. Seine Effizienz folgt aus der dünn besetzten
und lokalen Struktur des zugrunde liegenden Modells.
\end{summary}
